\chapter{Fontes, corpus e método}

\section{A imprensa como fonte e como objeto}

Tomo a imprensa periódica como fonte principal para investigar como a Caixa de Conversão foi apresentada, discutida e avaliada entre 1906 e 1914. Meu objetivo não é extrair dos jornais uma narrativa transparente dos acontecimentos, nem tratá-los apenas como repositórios de informações factuais. Interessa compreender a própria organização pública do debate: quais problemas ganharam visibilidade, quais atores receberam voz, que alternativas foram consideradas legítimas e quais vocabulários ligaram câmbio, moeda, café, crédito, dívida e interesses sociais.

A historiografia sobre a imprensa oferece uma advertência inicial importante. A relativa desconfiança que cercou o uso de jornais por historiadores brasileiros durante parte do século XX não decorria da ausência de informações, mas justamente de sua abundância, de sua proximidade com os acontecimentos e de seu envolvimento direto nas disputas políticas. Luca reconstrói essa mudança de estatuto e mostra que o uso histórico dos periódicos exige atenção simultânea ao conteúdo e às condições de produção, circulação e organização material do impresso \cite{luca2005}. Camargo, por sua vez, chama atenção para o risco de localizar no jornal apenas as confirmações desejadas pelo pesquisador, sobretudo quando uma palavra, uma passagem ou uma matéria é retirada da série e da conjuntura que lhe dão sentido \cite{camargo1969}. O problema não se resolve com a recusa da fonte, mas com um procedimento que preserve contexto, contraste e rastreabilidade.

Lapuente sistematiza exigências convergentes: examinar a materialidade do periódico, sua divisão interna, seus colaboradores, suas condições econômicas e políticas, além de fatores externos capazes de afetar aquilo que se publica ou omite \cite{lapuente2015}. No caso desta pesquisa, tais cuidados são especialmente relevantes. Os jornais estudados não eram observadores externos da política econômica. Eram empresas, atores políticos e espaços nos quais ministros, parlamentares, associações, banqueiros, correspondentes e articulistas procuravam intervir no debate. Uma mesma edição podia combinar editorial próprio, notícia de agência, discurso parlamentar, carta de banqueiro, anúncio e cotação cambial.

A imprensa da Primeira República também se formou num momento de transição empresarial e política. A expansão da grande imprensa alterou as relações entre jornal, publicidade, leitores e poder. Sodré descreve um campo marcado tanto por órgãos de oposição combativa quanto por jornais próximos aos governos, num ambiente em que a opinião publicada possuía valor político e econômico \cite{sodre1966}. Essa caracterização ajuda a localizar os periódicos, mas não funciona aqui como rótulo suficiente para interpretar qualquer matéria. Dizer que um jornal foi governista, oposicionista ou ligado a determinado partido oferece uma hipótese contextual, não uma prova antecipada de sua posição sobre cada aspecto da Caixa.

O princípio metodológico adotado pode ser resumido da seguinte forma: parcialidade, vínculos institucionais e organização material não são ruídos externos que precisam ser eliminados para alcançar uma notícia pura. São propriedades históricas da fonte e, portanto, parte do objeto. A questão deixa de ser se o jornal é parcial e passa a ser como essa parcialidade se constrói, em quais gêneros se manifesta, que vozes incorpora e como muda diante de conjunturas sucessivas.

\section{O corpus principal}

O corpus principal é composto por quatro diários preservados na Hemeroteca Digital Brasileira: \textit{O Paiz}, \textit{Correio da Manhã}, \textit{Correio Paulistano} e \textit{Gazeta de Notícias}. O recorte cobre 1906 a 1914, da criação da Caixa de Conversão à suspensão de suas operações de troco. Há cobertura identificada para todos esses anos nos três primeiros periódicos. A \textit{Gazeta de Notícias} não possui objetos recuperados para 1913 no censo atual. Essa ausência é mantida como ausência, sem imputação e sem substituir o jornal por outro veículo.

O conjunto corresponde a um censo do acervo digital que pôde ser identificado e recuperado na Biblioteca Nacional segundo os critérios documentados pelo projeto. Não corresponde a um censo de todas as edições historicamente publicadas, preservadas em qualquer instituição ou efetivamente lidas pelos contemporâneos. Essa distinção é substantiva: lacunas de preservação, catalogação e acesso digital condicionam o universo observável e devem acompanhar qualquer afirmação de cobertura.

Os quatro títulos foram escolhidos por combinarem importância nas praças do Rio de Janeiro e de São Paulo, continuidade suficiente no recorte e posições institucionais distintas. Segundo a caracterização de Sodré, \textit{O Paiz} se notabilizou pela proximidade com governos republicanos, ao passo que o \textit{Correio da Manhã}, fundado em 1901 por Edmundo Bittencourt, construiu sua identidade pública em torno da oposição \cite{sodre1966}. O \textit{Correio Paulistano}, fundado em 1854, tornou-se órgão do Partido Republicano Paulista e constitui uma fonte privilegiada para observar a defesa governista paulista e a circulação de argumentos ligados ao café. A \textit{Gazeta de Notícias}, fundada em 1875, integrou a modernização comercial da imprensa carioca e combinou ampla circulação de notícias, colaboração literária e intervenção política.

Essas localizações serão refinadas no texto empírico por meio da leitura das próprias peças. Elas não autorizam transformar todo conteúdo publicado em voz do periódico. Também não tornam o conjunto uma amostra representativa da imprensa brasileira em sentido estatístico. O valor comparativo do corpus está na possibilidade de acompanhar continuamente quatro instituições jornalísticas relevantes diante dos mesmos acontecimentos e de observar como suas práticas de seleção e enquadramento variaram no tempo.

O \textit{Estado de S. Paulo}, que utilizei na minha monografia de graduação sobre 1905 e 1906, não integra o censo principal desta dissertação. Quando aparecer, terá estatuto explícito de fonte externa ou de material reaproveitado daquele estudo, nunca de quinto componente do corpus. O mesmo vale para o diário do \textit{Jornal do Commercio}, cuja relevância histórica é reconhecida, mas cuja rota de acesso digital não permitiu sua inclusão serial nas mesmas condições dos quatro jornais.

\section{Fontes acessórias e hierarquia documental}

Os jornais são o centro da investigação, mas não podem ser interpretados isoladamente. Um segundo conjunto documental reúne debates e documentos do Congresso, legislação, mensagens presidenciais, relatórios ministeriais, cartas, telegramas financeiros e retrospectos comerciais. Essas fontes cumprem três funções. Primeiro, ajudam a reconstruir o acontecimento ao qual o jornal respondia. Segundo, permitem verificar se uma passagem publicada reproduz um documento ou se contém elaboração editorial própria. Terceiro, esclarecem o vocabulário técnico, as competências institucionais e os interesses em disputa.

Dois volumes de \textit{Documentos Parlamentares}, com 1.152 páginas em camada textual corrigida, oferecem uma base para confrontar a tramitação legislativa com sua apresentação na imprensa. As leis de 1906, 1910 e 1914 fixam os contornos jurídicos da criação, da mudança de paridade e da suspensão \cite{brasil1906caixa,brasil1910taxa16,brasil1914suspensao}. O \textit{Retrospecto Commercial} do \textit{Jornal do Commercio} está disponível em nove edições anuais, uma para cada ano do recorte. Ele serve como balanço da conjuntura e preserva correspondência de agentes financeiros, mas permanece fonte acessória, não substituto do diário ausente do censo.

Cartas e telegramas, inclusive os atribuídos aos Rothschild e a outros banqueiros estrangeiros, serão tratados simultaneamente como documentos e como atos de publicação. Seu conteúdo pode esclarecer a relação entre governo, credores e praças financeiras, enquanto o lugar que recebem na página revela como o jornal organiza a autoridade econômica. A transcrição de uma carta não equivale, por si, a endosso. A repetição, o título, o comentário introdutório, a posição na edição e eventuais respostas são elementos necessários para avaliar apropriação editorial.

Essa hierarquia me impede de dispersar a pesquisa num estudo geral de todas as fontes disponíveis. Os documentos acessórios ampliam a compreensão e qualificam a escrita, mas a pergunta principal continua sendo como os quatro jornais construíram o debate da Caixa. Quando uma conclusão depender principalmente de uma fonte parlamentar ou financeira, isso será indicado no texto.

\section{Construção e descrição do acervo}

O inventário digital contém 11.960 objetos: 3.087 de \textit{O Paiz}, 3.240 do \textit{Correio da Manhã}, 3.128 do \textit{Correio Paulistano} e 2.505 da \textit{Gazeta de Notícias}. Esses objetos somam 117.705 páginas, das quais 117.703 possuem texto extraído e duas foram registradas como vazias. A camada textual integral contém aproximadamente 4,32 bilhões de caracteres. Os números descrevem o estado auditado da base e serão atualizados no apêndice se o acervo for corrigido antes da entrega.

A busca tolerante às variações do OCR identificou 8.331 páginas com menção à Caixa de Conversão. Esse número não mede debate. A expressão aparece em cotações, balanços, boletins bancários, listas, referências incidentais e notícias nas quais a instituição não é objeto de controvérsia. A alta frequência nominal é, portanto, ponto de partida para recuperação documental, não denominador imediato de uma análise de posicionamento.

Uma camada amostral de catalogação tornou essa diferença visível. Entre 468 peças preservadas para exame, 453 receberam classificação humana quanto ao registro documental e quinze permanecem sem preenchimento. Das 453 classificadas, 239 foram consideradas substantivas, 169 operacionais ou rotineiras e 45 incidentais. Na caracterização provisória de forma textual, as 239 peças substantivas se distribuem em 117 notícias, 64 artigos, 28 editoriais, 20 telegramas e dez itens de outros tipos. Como parte dos rótulos de forma foi sugerida por modelo e ainda possui erro não medido, tais contagens descrevem a base de trabalho, não uma distribuição definitiva dos gêneros no universo.

O inventário se organiza, assim, em camadas que não devem ser confundidas: objetos digitais, páginas com texto, páginas com menção nominal, peças recuperadas e peças interpretadas como substantivas. Cada estreitamento responde a uma pergunta diferente. A ligação entre as camadas preserva o identificador do objeto, a página e, quando disponível, a localização da peça, permitindo retornar à imagem da Hemeroteca.

\section{Gêneros, vozes e atos de publicação}

A unidade mais útil para a leitura é a peça jornalística contextualizada na edição. Ela pode ser um editorial próprio, artigo assinado, notícia, seção livre ou espaço pago, discurso parlamentar reproduzido, carta ou telegrama transcrito, boletim operacional ou menção incidental. A categoria formal orienta a interpretação, mas não a decide sozinha. Um texto sem assinatura pode reproduzir despacho de agência; um artigo assinado pode ser escrito por integrante da direção; uma seção livre pode hospedar campanha política persistente; uma notícia pode receber título e comentário que marcam uma intervenção do jornal.

Por isso, distingo ao menos três níveis. O primeiro é a voz enunciadora, isto é, quem formula a afirmação. O segundo é o ato de publicação, que inclui a decisão do periódico de selecionar, titular, ordenar ou repetir o conteúdo. O terceiro é a apropriação editorial, que ocorre quando há evidência de que o jornal incorpora, comenta ou sustenta a posição publicada. Hospedar uma voz divergente é um dado sobre o periódico, mas não é automaticamente sua opinião.

Essa distinção é central para os documentos de terceiros. Discursos de David Campista, Leopoldo de Bulhões ou parlamentares podem circular em mais de um jornal com o mesmo núcleo textual. Telegramas de agências e correspondentes também podem ser compartilhados. A comparação de uma mesma janela cronológica entre os quatro periódicos permite observar tanto o material comum quanto as diferenças de seleção, destaque e comentário. Atribuições de posicionamento serão mais firmes quando apoiadas por editoriais, comentários próprios, recorrência coerente ou apropriação explícita. Nos demais casos, o texto falará em voz publicada, argumento hospedado ou espaço concedido.

\section{Procedimento de recuperação, catalogação e leitura}

A escala do acervo tornou necessário combinar procedimentos computacionais e leitura histórica. O fluxo começa pela preservação dos PDFs da Biblioteca Nacional, seus metadados e sua proveniência. Em seguida, utiliza o texto OCR embutido nos arquivos para formar uma camada pesquisável. Regras tolerantes a ruído recuperam menções e candidatos, sempre mantendo a ligação com o documento original. A partir daí, peças são segmentadas e catalogadas segundo forma, voz, atores, objetos de política, argumentos, vocabulário, marco cronológico e evidência textual.

\subsection{O instrumento computacional e suas guardas}

Modelos de linguagem entram nesta pesquisa em duas funções distintas, ambas de recuperação e descrição, nenhuma de julgamento sobre posição editorial. A primeira é a catalogação de janelas de texto, que organiza metadados de um conjunto extenso. A segunda é a conferência de transcrição contra a imagem do scan, onde o OCR falha. Descrevo as duas com a especificação que permite auditá-las, porque um instrumento de medição não declarado não é verificável.

A catalogação usa um prompt versionado, `catalogacao\_debates.md`, versão 0.1.0, com \textit{hash} de congelamento `61030ebfd7b93485`. Ele fixa vocabulário controlado para objeto de política, voz e agentes, proíbe completar com conhecimento externo e trata lista vazia como resposta válida e frequente. Mudança no prompt exige versão nova e registro datado, pela mesma razão que se registra a troca de um instrumento de laboratório. No piloto de julho de 2026, 48 janelas de seis mil caracteres, distribuídas pelos quatro jornais e pelas quatro fases, foram catalogadas duas vezes, por dois anotadores diferentes recebendo o mesmo prompt. O primeiro produziu 98 registros aceitos, o segundo 83.

A guarda central é a conferência mecânica de citação. Todo trecho que um anotador devolve como citação precisa ser, depois de normalizado, substring literal da janela de origem. A linha que não casa é rejeitada e nunca corrigida, porque corrigir a citação de um anotador destrói justamente a informação que a conferência produz. A taxa de rejeição por anotador é métrica de qualidade do lote e é publicada: 6,7 por cento e 5,7 por cento no piloto.

A conferência mecânica responde se a linha existe no OCR. Ela não responde o que está impresso na página, e essa diferença não é acadêmica: uma citação aprovada pode vir do OCR como \textit{111:1 lelcgraiimia de aos-sos banqueiros em Londres}, literal e inutilizável. Por isso há uma segunda camada, que lê a imagem do scan com o modelo `claude-sonnet-5`, sob o protocolo `conferencia-citacao-imagem/claude-sonnet-5 1.0.0`, sem raciocínio estendido, pedindo transcrição verbatim trecho a trecho. Cada página é lida duas vezes, em chamadas separadas, e cada leitura é ancorada no OCR que diz corrigir antes de ser comparada com a outra. A ancoragem vem primeiro de propósito: duas leituras idênticas e ambas erradas concordam perfeitamente, e a concordância entre elas não pode absolver o par de ter saído da passagem. Os vereditos são estável, instável e divergente do OCR. Na primeira aplicação, sobre 36 citações, 23 saíram estáveis, 11 instáveis e duas divergentes do OCR.

Duas ressalvas fecham o desenho, e as duas são limitações, não formalidades. Duas leituras do mesmo modelo não são dois anotadores independentes, e erro correlacionado atravessa as duas. Há caso medido em que as duas leituras omitiram a mesma palavra do original, justamente a palavra decisiva da passagem. Nenhum veredito desta camada, portanto, dispensa a leitura da página, e o que ela entrega é triagem, isto é, trocar conferir trinta e seis citações na imagem por conferir treze primeiro.

Meu argumento principal também não depende de uma pontuação holística que ordene cada edição entre ortodoxia e expansionismo. Tal escala pode ser examinada como instrumento auxiliar, mas comprime dimensões que variam separadamente e pode atribuir ao jornal uma posição presente apenas em discurso reproduzido. Concordância entre dois modelos é análise de sensibilidade descritiva, não validação humana e não validade de construto.

\lacuna{O relatório do piloto de catalogação registra as versões de interface dos dois anotadores, 2.1.220 e 0.145.0, mas não registra o identificador exato do modelo de cada um nem a temperatura de amostragem. A regra do projeto exige serviço, modelo, versão exata e prompt no \textit{output} de toda anotação. Recuperar esses campos dos registros da rodada, ou declarar a rodada como não reprodutível nesse aspecto, antes que qualquer contagem derivada dela entre no texto.}

O uso quantitativo da base será prioritariamente inventarial e descritivo. Ele permitirá apresentar cobertura por jornal e ano, composição documental, ocorrência de gêneros, distribuição de atores e saliência de objetos de debate. Esses resultados servem para localizar padrões e exceções que exigem leitura, não para substituir a interpretação. O enriquecimento dos metadados é valioso precisamente porque permite retornar de uma contagem a peças identificáveis e comparar diferentes dimensões sem reduzir todo o debate a um único número.

A análise histórica seguirá seis movimentos. Primeiro, reconstruir o acontecimento e a decisão institucional com apoio da legislação e da historiografia. Segundo, delimitar uma janela cronológica comum aos jornais. Terceiro, inventariar as peças recuperadas, seus gêneros e suas vozes. Quarto, identificar os objetos concretos de desacordo, os atores e os termos usados. Quinto, comparar seleção, destaque, comentário e mudança temporal entre periódicos. Sexto, confrontar o que aparece nas fontes com categorias da literatura, como ortodoxia, expansionismo, metalismo, defesa do café, credibilidade e disciplina monetária. As categorias historiográficas funcionarão como perguntas e hipóteses comparativas, não como etiquetas impostas de antemão.

\section{Limites e controle da interpretação}

O primeiro limite é a qualidade desigual do OCR. Uma medida de ruído aplicada ao corpus variou, no caso de \textit{O Paiz}, de 4,84 por cento em 1909 a 14,33 por cento em 1906. Diferenças desse porte tornam arriscadas comparações lexicais diretas entre jornais e anos. Ortografia histórica, colunas fragmentadas, páginas danificadas e erros de segmentação podem produzir tanto falsos negativos quanto formas artificiais de uma palavra. Por isso, frequência lexical será interpretada junto à qualidade documental e passagens decisivas serão conferidas na imagem.

O segundo limite é a relação entre silêncio e ausência. A falta de uma menção pode refletir desinteresse editorial, erro de OCR, lacuna do acervo ou escolha vocabular diferente. Em particular, a falta da \textit{Gazeta de Notícias} em 1913 impede uma comparação completa naquele ano. Não preencherei essa célula com inferência retrospectiva.

O terceiro limite é a desigualdade entre presença documental e posição editorial. O pequeno número de editoriais já identificado em algumas fases pode resultar da prática jornalística, do desenho da amostra ou de erro na identificação do gênero. Por isso, generalizações sobre cada periódico serão graduadas segundo a natureza da evidência e testadas contra artigos, notícias comentadas e episódios comparáveis.

Finalmente, este capítulo descreve um método de pesquisa em andamento. As contagens são auditáveis, mas a interpretação das peças continua sendo trabalho histórico. O objetivo inventarial é produzir a descrição mais exaustiva possível dentro das restrições do acervo e das ferramentas, deixando visíveis tanto o caminho entre página e argumento quanto as áreas em que a evidência ainda é insuficiente.

\lacuna{Acrescentar, antes da versão de qualificação, uma tabela sintética com os identificadores bibliográficos da Hemeroteca, a cobertura anual e as datas de auditoria de cada camada.}
